\GalSourcePageBoundary{078}{L0077}{49}{49}

Il \emph{faut} donc et il \emph{suffit} que l'équation qui donne cette fonction
des racines admette, quel que soit $X$, une valeur rationnelle.

Si l'équation proposée a tous ses coefficients rationnels, l'équation
auxiliaire qui donne cette fonction les aura tous aussi, et il
suffira de reconnaître si cette équation auxiliaire du degré
$1\mathbin{.}2\mathbin{.}3\ldots(n-2)$ a ou non une racine rationnelle, ce que l'on
sait faire.

C'est là le moyen qu'il faudrait employer dans la pratique. Mais
nous allons présenter le théorème sous une autre forme.

\begin{center}{\bfseries PROPOSITION VIII.}\end{center}
\GalSourceErrorMark{W09-SE006}

\noindent\textsc{Théorème.} --- {\itshape Pour qu'une équation irréductible de degré
premier soit soluble par radicaux, il \emph{faut} et il \emph{suffit} que deux
quelconques des racines étant connues, les autres s'en déduisent
rationnellement.}

Premièrement, il le faut, car la substitution
\[
x_{k},\qquad x_{ak+b}
\]
ne laissant jamais deux lettres à la même place, il est clair qu'en
adjoignant deux racines à l'équation, par la proposition IV, son
groupe devra se réduire à une seule permutation.

En second lieu, cela suffit; car, dans ce cas, aucune substitution
du groupe ne laissera deux lettres aux mêmes places. Par conséquent,
le groupe contiendra tout au plus $n(n-1)$ permutations.
Donc, il ne contiendra qu'une seule substitution circulaire (sans
quoi il y aurait au moins $n^{2}$ permutations). Donc, toute substitution
du groupe, $x_{k},x_{f(k)}$, devra satisfaire à la condition
\[
f(k+c)=f(k)+C.
\]
Donc, etc.

Le théorème est donc démontré.

\medskip
\noindent\hfill \textsc{E. G.}\hspace{3em}
